\documentclass{article}
\usepackage[utf8]{inputenc}
\usepackage[T1]{fontenc}
\usepackage[margin=2cm]{geometry}
\usepackage{hyperref}
\sloppy
\begin{document}
\section*{Historical source}
Encyclopaedia Britannica scan, source row 713. Parse 3.5 extracted the text. Super 3.5 VL corrected it from an enlarged crop of each block. It can still contain recognition errors.
Page headers, footers, and page numbers are removed.

\subsection*{Untitled continuation (E0062)}
\textit{Continues an entry from an earlier page.}

O V I ( 447 ) O U S

1. Aries, or ram and sheep, the horns of which are flapped like a half moon, and compressed. This animal is perhaps the most gentle and inoffensive, and at the same time the most timid and stupid, of all quadrupeds. In a flock of sheep, there is always a leader, whose steps the rest implicitly and blindly follow. If he runs over a precipice, the whole flock follow his example, and cannot be restrained though evident destruction be the consequence. The slightest noise, if uncommon, makes them run precipitantly against each other, without knowing the cause of their danger, or being able to extricate themselves. In snow, they remain fixed in the same spot, and continue obstinately in that situation, unless they are forced to move by the shepherd or his dog. But though dullardly in itself, though almost devoid of sentiment and mental qualities, this animal is of the most extensive utility to man. It affords us both food and clothing, besides the advantages we derive from its milk, its fleece, and its tallow.

Love, the liveliest sentiment in all animals, is the only passion which inspires the ram with any degree of vivacity; when under the influence of this passion, he becomes petulant, runs against his neighbours, and sometimes even attacks the shepherd himself. But the ewe, or female, though in the same situation, does not appear to be in the least affected; she seems to have no other instinct than that of simply admitting the embraces of the ram, eating her food, and recognizing her offspring; the lamb, on the other hand, is endowed with the same instinct of distinguishing its mother from the rest of the flock.

This animal, so soft and so simple in its nature and disposition, is likewise extremely weak and feeble in its condition: They can endure but little fatigue; whenever they run, their hearts palpitate, and their wind fails them; they are equally incapable of bearing heat or cold, snow or rain: They are subject to many diseases, most of which are contagious: They bring forth their young with pain and difficulty, and require more care and attention than any other domestic animal.

The ram is capable of generation at the age of 18 months; and the ewe can be impregnated when a year old. One ram is sufficient for 30 or 40 ewes. He ought to be strong, well proportioned; his head should be thick and strong, his front large, his eyes black, his nose flat, his neck thick, his body long and tall, his testicles mally, and his tail long. White is the best colour for a ram. The ewes whose wool is most plentiful, bushy, long, soft, and white, are most proper for breeders, especially when at the same time they are of a large size, have a thick neck, and move nimbly.

Ewes are generally in season from the end of October to the beginning of April. However, they may be impregnated at any time, if they be fed with stimulating nourishment, as fat water, and bread made of hemp-feed. They go with young about five months, and generally bring forth but one at a time, though frequently two. In warm climates, they may bring forth twice in a year; but in Britain, France, and most parts of Europe, only once. They give milk plentifully for seven or eight months: Their milk affords nourishment for children and country people; good cheese is made of it, especially when mixed with cow-milk. They live from 10 to 12 years, and are capable of bringing forth as long as they live.

When the male lambs are not intended to be kept for propagation, but fattened for food, they ought to be castrated at the age of five or six months. This operation is performed two ways: in the one, an incision is made and the testicles taken out; in the other, a ligature is tied tight round the scrotum, above the testicles, which soon destroys the vessels which nourish them. After castration, they are called wethers.

The ram, ewe, and weather, when one year old, lose the two foreteeth of the under jaw; fix months afterwards, they lose the two foreteeth next to these; and at the age of three years, the teeth are all replaced. The age of a ram may likewise be discovered by their horns, which always appear the first year, and frequently as soon as they are brought forth. These horns uniformly acquire an additional ring every year as long as the creature lives. The ewes commonly have no horns, but a kind of long protuberances in place of them: However, some of them have two, and some four horns.

In Spain, and the southern parts of Europe, the flocks are kept in folds or stables during the night; but in Britain, where there is now no danger from wolves, they are allowed to remain without, both night and day; which makes the animals more healthy, and their flesh a more whole food. Dry and mountainous grounds, where thyme and sheep's fescue grass abound, are the best for the pasturing sheep.

2. The Guineensis, or Guinea sheep, has pendulous ears, lax hairy dew lips, and a prominence on the hind part of the head. The wool is short, like that of a goat. It is a native of Guinea.

3. The firepliceros, or Cretan sleep, has frait cari- nated horns, twilted in a spiral manner. It is a native of Mount Iola.

\subsection*{OULNEY (E0063)}
OULNEY, a market-town of Buckinghamshire, situated nine miles fourth-east of Northampton.

\subsection*{OULZ (E0064)}
OULZ, a town of Italy, in the province of Piedmont, situated in E. long. 6° 30', N. lat. 45°.

\subsection*{OUNCE (E0065)}
OUNCE, a little weight, the fifteenth part of a pound avoirdupois, and the twelfth part of a pound troy.

cancer, in azobby. See Leo.

\subsection*{OUNOLE (E0066)}
OUNOLE, a market-town of Northamptonshire, situated on the river Nen, twenty-two miles north-east of Northampton.

\subsection*{OVOLLO (E0067)}
OVOLLO, or Ovum, in architecture, a round moulding, whose profile, or sweep, in the ionic and composite capitals, is usually a quadrant of a circle: whence it is also commonly called the quarter-round.

It is usually cut with the representation of eggs and an-
chors or arrows heads placed alternately.

Ouss, a river, which, riding in the north of Yorkshire, runs fourth-rail by York; and, combining that course, falls into the Trent.

\subsection*{OUSTLAGH (E0068)}
OUSTLAGH, or Ostriach Country, is a part of Asiatic Russia, extending along the river Irtis to its confluence with the river Obi, and from thence northward along the banks of the Obi and Jenisei, into the gulph of the Man- gusea, or the frozen ocean; and extending also along


\section*{Source record}
\begin{description}
\item[Scan ID] \texttt{britannica:309f765b1ecc83e0a7f7762e44b4a927a7555443:train:713}
\item[Image file] \texttt{page-0713-3db0c66f81f4.jpg}
\item[Image SHA-256] \texttt{3db0c66f81f48bd0948a7644d3f0c1c4213d3d49e54c500a317bf147695b1fc5}
\item[Dataset] NationalLibraryOfScotland/encyclopaedia\_britannica\_illustrated at revision \texttt{309f765b1ecc83e0a7f7762e44b4a927a7555443} (CC0-1.0)
\item[Extraction] model \texttt{nvidia/nemotron-3.5-parse}, raw reply \texttt{extract/raw/page-0713-3db0c66f81f4.json}
\item[Raw reply SHA-256] \texttt{177d6bd4d1df0b084cbb8f82ba4e25052152d590d581fcd1910aea5c9a1eb224}
\item[Correction] model \texttt{nvidia/Nemotron-3.5-Super-VL}, block crops, raw replies \texttt{correct/raw/crop\_nothink/page-0713-3db0c66f81f4-b*.json}; parse\_fallback: 2, super\_vl: 18
\item[Running header] none
\end{description}
\end{document}
